\documentclass[12pt]{article}
\usepackage{amsmath, amssymb, booktabs, geometry}
\geometry{margin=1in}
\setlength{\parindent}{0pt}
\setlength{\parskip}{0.8em}

\title{ECON 101: Macroeconomics \\ Solutions -- Quiz 3}
\author{Instructor: Muhebullah Karimzada}
\date{October 22,  2025}

\begin{document}
\maketitle

\hrule
\vspace{1em}

\section*{Part A – Multiple Choice Solutions (1 mark each)}

\textbf{1. CPI inflation from 152 to 160} \\
Inflation $= \dfrac{160 - 152}{152}\times 100 = \dfrac{8}{152}\times 100 \approx 5.26\% \Rightarrow 5.3\%$.\\
\textbf{Answer: (B)}.

---

\textbf{2. Fisher approximation} \\
Real interest $\approx i - \pi = 5\% - 3\% = 2\%$. (Exact Fisher: $\dfrac{1.05}{1.03}-1 \approx 1.94\%$.) \\
\textbf{Answer: (B)}.

---

\textbf{3. CPI coverage} \\
The CPI covers the prices of goods and services purchased by consumers, including imports, utilities, restaurant meals, \emph{and} used vehicles. With the options given, all listed items are typically included in the CPI. \\
\emph{Grading note:} Since no ``None of the above'' option was provided, accept that \textbf{all options are included in CPI}. If a single key is required, flag this item for revision (replace with a clearly excluded item such as ``corporate bond purchases''). 

---
\newpage
\textbf{4. Catch-up time: growth differential} \\
Solve for $t$ in $28{,}000(1.04)^t = 42{,}000(1.02)^t$.\\
$\displaystyle \left(\frac{1.04}{1.02}\right)^t = \frac{42{,}000}{28{,}000} = 1.5$.\\
$t = \dfrac{\ln(1.5)}{\ln(1.04/1.02)} \approx \dfrac{0.405465}{0.019414} \approx 20.9$ years $\approx 21$ years.\\
\textbf{Answer: (A)}.

---

\textbf{5. Early mover in recessions} \\
Investment is the most volatile component of spending and typically falls first. \\
\textbf{Answer: (D)}.

\vspace{0.5em}
\hrule
\vspace{0.5em}

\section*{Part B – Long Problem Solutions (15 marks)}

\textbf{6. Economic Growth and Inflation Analysis}

\textit{Data:}\\
\begin{tabular}{lcccc}
\toprule
Item & Qty (Y1) & Qty (Y2) & Price (Y1) & Price (Y2) \\
\midrule
Smartphones & 200 & 250 & \$500 & \$520 \\
Tablets     & 100 & 120 & \$800 & \$880 \\
Headphones  & 300 & 330 & \$100 & \$90 \\
\bottomrule
\end{tabular}

\subsection*{(a) Nominal GDP, Real GDP (Y2 in Y1 prices), and Real Growth}

\underline{Step 1: Nominal GDP (current-year prices)} \\
Year 1:
\[
\text{NGDP}_1 = (200)(500) + (100)(800) + (300)(100)
= 100{,}000 + 80{,}000 + 30{,}000 = \boxed{210{,}000}
\]
Year 2:
\[
\text{NGDP}_2 = (250)(520) + (120)(880) + (330)(90)
= 130{,}000 + 105{,}600 + 29{,}700 = \boxed{265{,}300}
\]

\underline{Step 2: Real GDP for Year 2 in Year 1 prices (base Y1)} \\
Hold prices at Y1 levels, update quantities to Y2:
\[
\text{RGDP}_{2|Y1\ \text{prices}} = (250)(500) + (120)(800) + (330)(100)
= 125{,}000 + 96{,}000 + 33{,}000 = \boxed{254{,}000}
\]
Note: $\text{RGDP}_1$ at Y1 prices equals $\text{NGDP}_1 = \boxed{210{,}000}$.

\underline{Step 3: Real GDP growth (base Y1)} \\
\[
\text{Real growth} = \frac{254{,}000 - 210{,}000}{210{,}000}\times 100
= \frac{44{,}000}{210{,}000}\times 100 \approx \boxed{20.95\% \ (\approx 21\%)}
\]

\subsection*{(b) GDP Deflator and Inflation}

\underline{Step 1: GDP Deflator} \\
\[
\text{Deflator}_t = \frac{\text{NGDP}_t}{\text{RGDP}_t}\times 100
\]
Year 1: $\text{Deflator}_1 = \dfrac{210{,}000}{210{,}000}\times 100 = \boxed{100.00}$. \\
Year 2: $\text{Deflator}_2 = \dfrac{265{,}300}{254{,}000}\times 100 \approx \boxed{104.45}.
$

\underline{Step 2: Inflation (deflator-based)} \\
\[
\pi \approx \frac{104.45 - 100.00}{100.00}\times 100 = \boxed{4.45\%}
\]

\subsection*{(c) Real GDP per Capita Growth}

Population: $N_1=800$, $N_2=820$. \\
\[
\text{RGDPpc}_1 = \frac{210{,}000}{800} = \boxed{262.50}
\quad\text{and}\quad
\text{RGDPpc}_2 = \frac{254{,}000}{820} \approx \boxed{309.76}
\]
\[
\text{Growth in RGDP per capita} = \frac{309.76 - 262.50}{262.50}\times 100 \approx \boxed{18.10\%}
\]

\subsection*{(d) Living Standards Assessment}

Real GDP grew by about $21\%$, population rose by $2.5\%$ (from 800 to 820), and \emph{real GDP per capita} increased by about $18.1\%$. Despite moderate inflation (about $4.45\%$ via the deflator), the substantial rise in \emph{real} output per person indicates that living standards \textbf{improved markedly}.

\vspace{0.5em}
\hrule
\vspace{0.5em}

\begin{center}
\textit{End of Instructor Solutions}
\end{center}

\end{document}
